\documentclass[12pt,a4paper]{article}

\usepackage[utf8]{inputenc}
\usepackage[T1]{fontenc}
\usepackage[romanian]{babel}
\usepackage[a4paper,margin=2.5cm]{geometry}
\usepackage{graphicx}
\usepackage{float}
\usepackage{xcolor}
\usepackage{listings}
\usepackage{booktabs}
\usepackage{longtable}
\usepackage{array}
\usepackage{tabularx}
\usepackage{enumitem}
\usepackage{titlesec}
\usepackage{fancyhdr}
\usepackage[hidelinks]{hyperref}
\usepackage{caption}

\pagestyle{fancy}
\fancyhf{}
\rhead{\small Deskly / AuthX -- Raport de securitate}
\lhead{\small Dezvoltarea Aplicatiilor Software Securizate}
\rfoot{\thepage}
\renewcommand{\headrulewidth}{0.4pt}

\titleformat{\section}{\Large\bfseries}{\thesection.}{0.6em}{}
\titleformat{\subsection}{\large\bfseries}{\thesubsection}{0.6em}{}

\definecolor{codebg}{RGB}{248,248,248}
\definecolor{codekey}{RGB}{0,0,180}
\definecolor{codestr}{RGB}{170,30,30}
\definecolor{codecom}{RGB}{40,120,40}
\definecolor{codenum}{RGB}{120,120,120}

\lstdefinestyle{py}{
  language=Python,
  basicstyle=\ttfamily\footnotesize,
  keywordstyle=\color{codekey}\bfseries,
  stringstyle=\color{codestr},
  commentstyle=\color{codecom}\itshape,
  numbers=left,
  numberstyle=\tiny\color{codenum},
  frame=single,
  backgroundcolor=\color{codebg},
  breaklines=true,
  breakatwhitespace=true,
  showstringspaces=false,
  tabsize=4,
  captionpos=b,
  xleftmargin=1em,
}
\lstdefinestyle{sh}{
  language=bash,
  basicstyle=\ttfamily\footnotesize,
  keywordstyle=\color{codekey}\bfseries,
  stringstyle=\color{codestr},
  commentstyle=\color{codecom}\itshape,
  frame=single,
  backgroundcolor=\color{codebg},
  breaklines=true,
  showstringspaces=false,
  captionpos=b,
  xleftmargin=1em,
}
\lstdefinestyle{sqltext}{
  language=SQL,
  basicstyle=\ttfamily\footnotesize,
  keywordstyle=\color{codekey}\bfseries,
  stringstyle=\color{codestr},
  commentstyle=\color{codecom}\itshape,
  frame=single,
  backgroundcolor=\color{codebg},
  breaklines=true,
  showstringspaces=false,
  captionpos=b,
  xleftmargin=1em,
}

\graphicspath{{./}}
\newcommand{\captura}[3]{
  \begin{figure}[H]
    \centering
    \IfFileExists{#1}{
      \includegraphics[width=0.85\textwidth,keepaspectratio]{#1}
    }{
      \fbox{\parbox{0.78\textwidth}{\centering\vspace{2.2cm}\textit{[Aici se insereaza captura: #1]}\vspace{2.2cm}}}
    }
    \caption{#2}\label{fig:#3}
  \end{figure}
}

\setlist{topsep=2pt,itemsep=2pt,parsep=0pt}

\begin{document}

\begin{titlepage}
\begin{center}
\vspace*{1.5cm}

{\large UNIVERSITATEA DIN BUCURESTI}\\[0.2cm]
{\large FACULTATEA DE MATEMATICA SI INFORMATICA}\\[0.2cm]
{\large Specializarea Informatica, Anul III, Semestrul II}\\[2.5cm]

{\Huge\bfseries Raport de securitate}\\[0.6cm]
{\LARGE\itshape Proiect 2 -- Break the Login}\\[0.4cm]
{\Large Atacarea si securizarea autentificarii}\\[0.6cm]
{\large Aplicatia ``Deskly / AuthX''}\\[3cm]

\begin{tabular}{rl}
\textbf{Student:}     & Nume Prenume \\
\textbf{Grupa:}       & NNN \\
\textbf{Username VM:} & \texttt{rares.stancu@deskly-vm} \\
\textbf{Email:}       & \texttt{stancurares10@gmail.com} \\
\textbf{Disciplina:}  & Dezvoltarea Aplicatiilor Software Securizate \\
\textbf{Data:}        & aprilie 2026 \\
\end{tabular}

\vfill
{\small Documentul este insotit de codul sursa (versiunea vulnerabila si versiunea securizata),
de scripturi de smoke-test (\texttt{smoke\_vuln.py}, \texttt{smoke\_fixed.py}),
de un clip video de 5--10 minute care
demonstreaza live atacurile si fix-urile.}
\end{center}
\end{titlepage}

\tableofcontents
\listoffigures
\listoftables
\newpage

\section*{Rezumat}
\addcontentsline{toc}{section}{Rezumat}

Lucrarea de fata documenteaza analiza de securitate si remedierea
vulnerabilitatilor unei aplicatii web de autentificare construita ca
parte a Proiectului~2 ``Break the Login''. Aplicatia, denumita generic
\textit{Deskly / AuthX}, expune in mod intentionat un mecanism de
autentificare slab construit (versiunea \textit{vulnerable}). Pe acelasi
nucleu am construit, intr-o a doua iteratie, o versiune securizata
(\textit{fixed}) care implementeaza controalele cerute de specificatie
plus controalele suplimentare semnalate in diagrama de arhitectura
(prevenirea SQL injection si a IDOR, audit logging si Content Security
Policy).

Raportul respecta structura impusa de baremul de evaluare: introducere,
setup-ul mediului de lucru, implementarea MVP, prezentarea tehnica a
vulnerabilitatilor, demonstrarea atacurilor (PoC) cu capturi de ecran,
analiza impactului, descrierea fix-urilor, re-testul si capitolul dedicat
audit-ului. Toate cele 6 vulnerabilitati obligatorii (4.1--4.6) sunt
documentate, plus alte 3 controale derivate din arhitectura (4.7 SQLi,
4.8 IDOR, 4.9 audit). Fiecare atac este reproductibil pas cu pas, iar
fix-ul corespunzator este re-testat cu acelasi PoC, demonstrand
inchiderea vulnerabilitatii.

\newpage

\section{Introducere}\label{sec:intro}

\subsection{Context si obiectiv}
Proiectul ``Break the Login'' urmareste consolidarea cunostintelor de
securitate aplicativa prin parcurgerea ciclului complet
\textbf{Build $\rightarrow$ Hack $\rightarrow$ Fix $\rightarrow$ Retest}
asupra unui sistem de autentificare. Studentul joaca alternativ rolul de
\textit{security tester} (identifica si demonstreaza vulnerabilitatile)
si rolul de \textit{developer} (implementeaza controalele corecte).

Conform specificatiei, accentul cade exclusiv pe:
\begin{itemize}
  \item logica de autentificare (Register, Login);
  \item gestionarea parolelor (stocare, complexitate);
  \item gestionarea sesiunilor sau a token-urilor;
  \item resetarea parolei.
\end{itemize}

Functionalitatile de business (crearea, vizualizarea si cautarea de
ticket-uri) sunt subordonate acestui obiectiv: ele permit demonstrarea
problemelor de control acces (IDOR) si a injectiilor SQL pe un endpoint
de cautare.

\subsection{Scenariu}
Compania fictiva \textit{AuthX} dezvolta o aplicatie interna numita
\textit{Deskly}, folosita de doua roluri: \texttt{ANALYST} si
\texttt{MANAGER}. Echipa de securitate a semnalat probleme grave la
componenta de autentificare. Studentul preia codul aflat ``in productie''
(versiunea vulnerabila), executa un mini-pentest si livreaza o versiune
securizata, plus prezentul raport.

\subsection{Arhitectura aplicatiei}
Aplicatia este una clasica de tip server-rendered: backend Flask
(Python~3.13), templating Jinja2, persistenta in SQLite, sesiuni
server-side (tabela \texttt{sessions}, cookie \texttt{session\_id}).
Structura pe nivele este urmatoarea:

\begin{itemize}
  \item \textbf{Web UI} -- formulare HTML pentru Register, Login,
        Forgot, Reset, plus paginile de tickets, search si audit.
  \item \textbf{Backend API} -- rutele Flask din \texttt{app.py}:
        autentificare, CRUD ticket, search, audit viewer (doar v2).
  \item \textbf{Security Controls (in v2)} -- input validation, output
        encoding, error handling, AuthZ (RBAC + ownership), CSRF,
        session/JWT mgmt, audit logging, CSP.
  \item \textbf{Data Layer} -- repository-uri (UserRepo, TicketRepo,
        AuditRepo) cu query-uri parametrizate.
  \item \textbf{Database} -- SQLite, trei tabele: \texttt{users},
        \texttt{tickets}, \texttt{audit\_logs}, plus tabela
        \texttt{sessions} pentru sesiuni server-side.
\end{itemize}

Diagrama de arhitectura furnizata in specificatie pune in evidenta
controalele suplimentare cerute in versiunea securizata: prevenirea
injectiei (prepared statements + input validation), prevenirea XSS
(output encoding in templates), tratarea exceptiilor fara stack-trace,
verificarea ownership-ului pe fiecare endpoint, protectie CSRF si
hardening de cookie-uri.

\subsection{Structura raportului}
Raportul urmeaza, in ordine, sectiunile cerute de barem (Setup mediu,
Implementare MVP, Prezentare vulnerabilitati, Demonstrare atac, Analiza
impact, Fix, Re-test, Audit \& logging, Concluzii). La final am inclus
un capitol dedicat criteriilor de bonus (CSP, rate limiting avansat,
CI/CD, teste automate de securitate).

\newpage
\section{Setup mediu}\label{sec:setup}

\subsection{Masina virtuala}
Aplicatia este rulata exclusiv pe o masina virtuala pentru a respecta
cerinta de izolare a mediului de testare. VM-ul foloseste Ubuntu~24.04
LTS, alocate 2~vCPU si 4~GB RAM, retea NAT cu port forwarding 5000.
Hostname-ul si username-ul sunt vizibile in fiecare captura de ecran,
asa cum cere baremul.

\begin{lstlisting}[style=sh, caption={Identificarea mediului de lucru}]
rares.stancu@deskly-vm:~$ whoami
rares.stancu
rares.stancu@deskly-vm:~$ hostname
deskly-vm
rares.stancu@deskly-vm:~$ uname -a
Linux deskly-vm 6.8.0-31-generic ...
rares.stancu@deskly-vm:~$ date
Wed Apr 29 11:42:01 UTC 2026
\end{lstlisting}

\captura{identificare.png}{Identificarea studentului si a VM-ului: \texttt{whoami}, \texttt{hostname}, \texttt{uname -a} si \texttt{date} ruleaza in acelasi terminal; bara VirtualBox confirma data sistemului}{setup-id}

\subsection{Stack tehnologic}
Tehnologiile alese si argumentarea lor (versiunile sunt fixate in
fisierele \texttt{requirements.txt} ale celor doua versiuni):
\begin{itemize}
  \item \textbf{Python 3.13} -- versiune curenta, suport bcrypt 5.0.
  \item \textbf{Flask 3.1} -- framework web minimal, suficient pentru
        scopul didactic; nu introduce protectii implicite (potrivit
        pentru a expune vulnerabilitatile in v1).
  \item \textbf{SQLite 3} -- DB locala, fara dependente externe;
        baremul accepta SQLite explicit. Permite usor demonstratiile
        de SQL injection (engine real, sintaxa standard).
  \item \textbf{bcrypt 5.0} -- folosit doar in v2 pentru hash-area
        parolelor; cost factor 12 (echilibru intre rezistenta si
        latenta sub 250~ms).
  \item \textbf{Jinja2} -- inclus in Flask, escapare HTML automata
        pentru prevenirea XSS.
\end{itemize}

\subsection{Structura proiectului}
Cele doua versiuni sunt fizic separate in directoarele
\texttt{vulnerable/} si \texttt{fixed/}, fiecare cu propriul
\texttt{app.py}, propriul director \texttt{templates/} si propriul
fisier de baza de date generat la prima rulare.

\begin{lstlisting}[style=sh, caption={Arborele proiectului}]
deskly/
|-- vulnerable/
|   |-- app.py                 
|   |-- requirements.txt
|   |-- templates/             
|   `-- deskly_vuln.db         
|-- fixed/
|   |-- app.py                 
|   |-- requirements.txt       
|   |-- templates/             
|   `-- deskly_fixed.db
|-- smoke_vuln.py              
|-- smoke_fixed.py             
\end{lstlisting}

\subsection{Pornirea aplicatiilor}
\begin{lstlisting}[style=sh, caption={Pornirea versiunii vulnerabile}]
rares.stancu@deskly-vm:~$ cd ~/proiect/vulnerable
rares.stancu@deskly-vm:~/proiect/vulnerable$ python3 -m venv .venv
rares.stancu@deskly-vm:~/proiect/vulnerable$ source .venv/bin/activate
(.venv) rares.stancu@deskly-vm:~/proiect/vulnerable$ pip install -r requirements.txt
(.venv) rares.stancu@deskly-vm:~/proiect/vulnerable$ python app.py
 * Running on http://127.0.0.1:5000
\end{lstlisting}

\captura{pornit.png}{Aplicatia Flask ascultand pe \texttt{127.0.0.1:5000}: in captura este pornita versiunea \texttt{fixed} (din \texttt{\textasciitilde/proiecte/deskly/fixed}); pentru \texttt{vulnerable} comenzile sunt identice, doar directorul difera}{setup-server}

Pentru v2 procedeul este identic, intr-un terminal separat, dupa ce v1
este oprit (ambele folosesc port-ul 5000).

\subsection{Conturi seed si initializare DB}
La prima pornire, fiecare versiune executa \texttt{init\_db()} care
creeaza schema (\texttt{users}, \texttt{tickets}, \texttt{audit\_logs},
\texttt{sessions}) si insereaza doua conturi de test, plus doua
ticket-uri. Conturile sunt:

\begin{table}[H]
\centering
\begin{tabular}{lll}
\toprule
\textbf{Email} & \textbf{Parola} & \textbf{Rol} \\
\midrule
\texttt{manager@deskly.local} & \texttt{Manager123!} & MANAGER \\
\texttt{analyst@deskly.local} & \texttt{Analyst123!} & ANALYST \\
\bottomrule
\end{tabular}
\caption{Conturi seed (identice in v1 si v2)}
\label{tab:seed}
\end{table}

In v1, hash-urile sunt MD5; in v2, bcrypt cost 12. Acelasi script de
initializare ruleaza pe ambele versiuni, dar ramura de inserare a
parolei difera intre cele doua \texttt{app.py}.

\newpage
\section{Implementarea MVP}\label{sec:mvp}

\subsection{Schema bazei de date}
Schema este identica intre v1 si v2 ca structura, ca sa permita
compararea atacurilor pe acelasi set de date. Diferentele apar la
continut (hash-ul parolei, prezenta sau absenta inregistrarilor de
audit).

\begin{lstlisting}[style=sqltext, caption={Schema bazei de date (extras din \texttt{init\_db()})}]
CREATE TABLE users (
    id INTEGER PRIMARY KEY AUTOINCREMENT,
    email TEXT UNIQUE NOT NULL,
    password_hash TEXT NOT NULL,
    role TEXT NOT NULL DEFAULT 'ANALYST',
    created_at TIMESTAMP DEFAULT CURRENT_TIMESTAMP,
    locked INTEGER NOT NULL DEFAULT 0
);
CREATE TABLE tickets (
    id INTEGER PRIMARY KEY AUTOINCREMENT,
    title TEXT NOT NULL,
    description TEXT,
    severity TEXT NOT NULL DEFAULT 'LOW',
    status TEXT NOT NULL DEFAULT 'OPEN',
    owner_id INTEGER NOT NULL,
    created_at TIMESTAMP DEFAULT CURRENT_TIMESTAMP,
    updated_at TIMESTAMP DEFAULT CURRENT_TIMESTAMP,
    FOREIGN KEY (owner_id) REFERENCES users(id)
);
CREATE TABLE audit_logs (
    id INTEGER PRIMARY KEY AUTOINCREMENT,
    user_id INTEGER, action TEXT, resource TEXT,
    resource_id TEXT, timestamp TIMESTAMP DEFAULT CURRENT_TIMESTAMP,
    ip_address TEXT
);
CREATE TABLE sessions (
    id TEXT PRIMARY KEY,
    user_id INTEGER NOT NULL,
    created_at TIMESTAMP DEFAULT CURRENT_TIMESTAMP,
    -- in v2 se adauga: expires_at, ip_address, user_agent, revoked
    FOREIGN KEY (user_id) REFERENCES users(id)
);
\end{lstlisting}

Cheia de prevenire IDOR este coloana \texttt{tickets.owner\_id}: in v2,
fiecare endpoint care returneaza un ticket verifica explicit
\texttt{ticket.owner\_id == current\_user.id} sau ca rolul este
\texttt{MANAGER}.

\subsection{Inregistrarea utilizatorilor (Register)}
Formularul colecteaza email-ul si parola, apoi trimite \texttt{POST
/register}. Versiunea vulnerabila accepta orice parola non-vida si
foloseste MD5 fara salt; versiunea securizata aplica
\texttt{PASSWORD\_POLICY\_RE} si bcrypt.

\begin{lstlisting}[style=py, caption={Register (extras din \texttt{vulnerable/app.py})}]
@app.route("/register", methods=["GET", "POST"])
def register():
    error = None
    if request.method == "POST":
        email = request.form.get("email", "").strip()
        password = request.form.get("password", "")
        confirm = request.form.get("confirm", "")
        if not email or not password:
            error = "Email and password required"
        elif password != confirm:
            error = "Passwords do not match"
        else:
            db = get_db()
            db.execute(
                "INSERT INTO users (email, password_hash, role) "
                "VALUES (?, ?, 'ANALYST')",
                (email, hash_password(password)),  
            )
            db.commit()
            return redirect(url_for("login"))
    return render_template("register.html", error=error)
\end{lstlisting}

\subsection{Autentificare si gestionarea sesiunii}
Sesiunile sunt server-side: la login se insereaza un rand in tabela
\texttt{sessions} si se returneaza un cookie \texttt{session\_id} care
contine doar ID-ul opac. Aceasta arhitectura permite invalidarea reala
a sesiunii la logout (in v2) -- in v1 logout-ul sterge doar cookie-ul
client.

\begin{lstlisting}[style=py, caption={Login si emiterea cookie-ului in v1}]
token = secrets.token_hex(16)
db.execute("INSERT INTO sessions (id, user_id) VALUES (?, ?)",
           (token, user["id"]))
db.commit()
resp = make_response(redirect(url_for("dashboard")))
resp.set_cookie(SESSION_COOKIE, token, max_age=30 * 24 * 3600)
return resp
\end{lstlisting}

\subsection{Resetarea parolei}
Fluxul ``Forgot password'' afiseaza tokenul direct in pagina (in mediul
real ar fi trimis prin email). In v1 tokenul are forma
\texttt{base64url(user\_id : timestamp)}, fara nicio inregistrare
server-side. In v2 tokenul este aleator (32~bytes), stocat hash-uit, cu
expirare 1~ora si invalidat dupa folosire.

\subsection{CRUD Tickets si Search}
Listarea ticketelor respecta rolul (un \texttt{ANALYST} vede doar ce a
creat, un \texttt{MANAGER} vede tot). Endpoint-ul \texttt{/ticket/<id>}
in v1 nu verifica ownership-ul, ceea ce conduce la IDOR. Endpoint-ul
\texttt{/ticket/search} in v1 concateneaza string SQL, ceea ce permite
SQL injection.

\begin{lstlisting}[style=py, caption={Search vulnerabil (concatenare SQL)}]
@app.route("/ticket/search")
@login_required
def ticket_search():
    q = request.args.get("q", "")
    sql = ("SELECT id, title, description, severity, status, owner_id "
           "FROM tickets WHERE title LIKE '
           "OR description LIKE '
    try:
        rows = get_db().execute(sql).fetchall()
    except sqlite3.Error as e:
        # Leaks SQL-engine errors -- ajuta atacatorul sa tuneze payload-ul.
        return render_template("ticket_search.html", q=q, tickets=[],
                               error=f"SQL error: {e}")
    return render_template("ticket_search.html", q=q, tickets=rows)
\end{lstlisting}

\newpage
\section{Prezentarea vulnerabilitatilor (mapare OWASP)}\label{sec:vuln}

Fiecare vulnerabilitate este mapata pe categoria corespunzatoare din
\textit{OWASP Top~10 (2021)} si pe controlul corespunzator din
\textit{OWASP ASVS v4.0.3}. Maparea face explicita relatia dintre
problemele identificate si standardele recunoscute industrial.

\begin{longtable}{p{0.04\textwidth} p{0.22\textwidth} p{0.30\textwidth} p{0.34\textwidth}}
\toprule
\textbf{ID} & \textbf{Vulnerabilitate} & \textbf{OWASP Top 10 (2021)} & \textbf{ASVS v4.0.3} \\
\midrule
\endhead
4.1 & Password policy slaba & A07: Identification and Authentication Failures & V2.1.1 (lungime min.~12), V2.1.7 (banned passwords) \\
4.2 & Stocare nesigura (MD5) & A02: Cryptographic Failures & V2.4.1 (Argon2/bcrypt), V2.4.4 (cost suficient) \\
4.3 & Brute-force / no rate limit & A07 + A04: Insecure Design & V2.2.1 (anti-automation), V11.1.5 (lockout) \\
4.4 & User enumeration & A01: Broken Access Control / A07 & V2.2.6 (mesaj generic), V2.10.4 (timp uniform) \\
4.5 & Sesiune nesigura & A07 + A02 & V3.4.1 (HttpOnly), V3.4.2 (SameSite), V3.3.1 (rotatie) \\
4.6 & Reset token nesigur & A07 + A04 & V2.5.1 (token aleator), V2.5.4 (expirare), V2.5.7 (one-time) \\
4.7 & SQL injection & A03: Injection & V5.3.4 (parameterized queries) \\
4.8 & IDOR & A01: Broken Access Control & V4.2.1 (object-level authZ) \\
4.9 & Lipsa audit & A09: Security Logging Failures & V7.1.1, V7.1.3 (logarea evenimentelor) \\
\bottomrule
\caption{Maparea vulnerabilitatilor pe OWASP Top 10 si ASVS}
\label{tab:owasp}
\end{longtable}

Pentru fiecare vulnerabilitate este descris in continuare:
(i)~mecanismul tehnic prin care apare in v1, (ii)~PoC-ul reproductibil,
(iii)~impactul asupra triadei CIA si asupra business-ului fictiv,
(iv)~fix-ul implementat in v2, cu trimitere la cod, (v)~re-testul.

\subsection{Sumarizare tehnica per vulnerabilitate}
Tabelul de mai jos este o vedere compacta inainte de a intra in
detaliile fiecarui atac. Coloana ``Loc cod v1'' indica functia in care
apare problema; coloana ``Fix v2'' indica zona modificata.

\begin{longtable}{p{0.04\textwidth} p{0.18\textwidth} p{0.34\textwidth} p{0.34\textwidth}}
\toprule
\textbf{ID} & \textbf{Pe scurt} & \textbf{Loc cod v1} & \textbf{Fix v2} \\
\midrule
4.1 & Parole triviale & \texttt{register()}, \texttt{reset()} & \texttt{PASSWORD\_POLICY\_RE} (8+, U, d, sym) \\
4.2 & MD5 fara salt & \texttt{hash\_password()} & \texttt{bcrypt.hashpw(gensalt(rounds=12))} \\
4.3 & Brute-force liber & \texttt{login()} & \texttt{failed\_attempts} + \texttt{lock\_until} (5 / 15min) \\
4.4 & Mesaje distincte & \texttt{login()}, \texttt{forgot()} & Mesaj generic + bcrypt dummy timing \\
4.5 & Cookie + zombie & \texttt{login()}, \texttt{logout()} & HttpOnly + SameSite + rotatie + revoke \\
4.6 & Token previzibil & \texttt{forgot()}, \texttt{reset()} & \texttt{secrets.token\_urlsafe} + hash + 1h \\
4.7 & Concat. SQL & \texttt{ticket\_search()} & Parametrizare cu \texttt{?} \\
4.8 & IDOR & \texttt{ticket\_detail()} & Verif. \texttt{owner\_id} sau MANAGER \\
4.9 & Fara audit & peste tot & Helper \texttt{audit()} pe toate ramurile \\
\bottomrule
\caption{Sumar vulnerabilitati cu pointer la cod}
\label{tab:vuln-summary}
\end{longtable}

\newpage
\section{Demonstrarea atacurilor (PoC)}\label{sec:poc}

Pentru fiecare atac sunt prezentate: (a) request-ul HTTP relevant,
(b) o captura de ecran cu rezultatul, (c) interpretarea raspunsului.
inclus la livrare.

\subsection{4.1 -- Acceptarea parolelor triviale}
\textbf{Pas 1.} Trimitere formular \texttt{POST /register} cu o parola
de un caracter:
\begin{lstlisting}[style=sh]
POST /register HTTP/1.1
Host: 127.0.0.1:5000
Content-Type: application/x-www-form-urlencoded

email=weak_rares@x.local&password=a&confirm=a
\end{lstlisting}

\textbf{Raspuns.} HTTP 302 catre \texttt{/login}. Niciun mesaj de
eroare; rand nou in tabela \texttt{users}.

\captura{vulnParola.png}{Dupa register cu parola de un caracter, contul \texttt{weak.@gmail.com} ajunge direct pe \texttt{/dashboard} cu rolul \texttt{ANALYST}; in dreapta, randul aferent stocat in \texttt{deskly\_vuln.db} (hash MD5 in coloana \texttt{password\_hash})}{poc-41}

\subsection{4.2 -- Recuperarea parolelor din DB}
Cu acces la fisierul \texttt{deskly\_vuln.db}, hash-urile sunt vizibile
direct in coloana \texttt{password\_hash}:
\begin{lstlisting}[style=sh]
$ sqlite3 deskly_vuln.db "SELECT email, password_hash FROM users;"
manager@deskly.local | b65a1b6aef38e1ed3d62c73cd3e63dac
analyst@deskly.local | 91d5a6f1f0c0d47c9c8d3d3f3a5d8a2a
$ python3 -c "import hashlib; print(hashlib.md5(b'Manager123!').hexdigest())"
b65a1b6aef38e1ed3d62c73cd3e63dac
\end{lstlisting}
\textbf{Confirmare.} Hash-ul recalculat e identic cu cel din DB.

\captura{parolaMd5.png}{Tabela \texttt{users} din \texttt{deskly\_vuln.db}: hash-urile au exact 32 caractere hex, semnatura tipica MD5 fara salt; sunt usor de cracat offline}{poc-42}

\subsection{4.3 -- Brute-force fara lockout}
Pentru a demonstra ca \texttt{POST /login} accepta un numar nelimitat de
incercari, am rulat un loop \texttt{curl} cu 6 parole succesive pe
contul \texttt{manager@deskly.local}. Primele 5 (\texttt{123456},
\texttt{password}, \texttt{qwerty}, \texttt{letmein}, \texttt{admin})
au primit toate raspuns HTTP \textbf{200} cu pagina de eroare, iar a
6-a incercare cu parola corecta \texttt{Manager123!} a reusit la prima
trecere -- raspuns \textbf{302} catre \texttt{/dashboard}. In tabela
\texttt{users}, coloana \texttt{locked} a ramas \texttt{0} pentru toti
utilizatorii: nu exista mecanism de lockout, captcha sau delay.

\captura{rateLimit.png}{Loop \texttt{curl} pe \texttt{/login}: 5 esecuri (status 200) urmate de un succes la a 6-a incercare (status 302); randul \texttt{manager@deskly.local} ramane \texttt{locked=0} in DB}{poc-43}

\subsection{4.4 -- Enumerare utilizatori}
Trimiterea unui email inexistent vs. unui email valid produce mesaje
distincte:
\begin{itemize}
  \item \texttt{ghost@nope.local} $\rightarrow$ ``User does not exist''
  \item \texttt{manager@deskly.local} $\rightarrow$ ``Incorrect password''
\end{itemize}
Pe langa text, lungimea raspunsului HTTP difera, ceea ce permite
automatizarea atacului prin scripturi \texttt{curl} pe o lista de
email-uri (filtrare dupa coloana \texttt{size}).

\captura{userIncorect.png}{Email inexistent (\texttt{ghost@nope.local}): mesajul \emph{``User does not exist''} il anunta pe atacator ca poate trece la urmatorul candidat}{poc-44a}

\captura{incorrectPass.png}{Email valid + parola gresita (\texttt{manager@deskly.local} / \texttt{wrong}): mesajul diferit \emph{``Incorrect password''} confirma existenta contului}{poc-44b}

\captura{userEnumeration.png}{Atac timing-based cu \texttt{enum\_timing.sh}: ramura ``user gasit'' calculeaza MD5 si raspunde mai lent (\(\sim\)10~ms vs.\ \(\sim\)2~ms pentru ramura ``user lipsa''), permitand enumerarea chiar daca textul ar fi identic}{poc-44c}

\subsection{4.5 -- Cookie nesigur si sesiune ``zombie''}
\textbf{Pas 1.} Inspectia cookie-ului in DevTools arata absenta
flag-urilor HttpOnly, Secure si SameSite, plus o expirare de 30 zile.
\textbf{Pas 2.} Apelul \texttt{document.cookie} din Console returneaza
valoarea cookie-ului -- proba ca un payload XSS l-ar putea exfiltra.
\textbf{Pas 3.} Dupa \texttt{Logout}, replay cu \texttt{curl} pe
\texttt{/dashboard} folosind cookie-ul capturat reuseste:
\begin{lstlisting}[style=sh]
$ curl -s -b "session_id=8c31fa..." http://127.0.0.1:5000/dashboard \
       | grep -o "<title>.*</title>"
<title>Dashboard - Deskly</title>
\end{lstlisting}

\captura{document.cookie.png}{Cookie \texttt{session\_id} in v1: in browser-ul logat ca \texttt{manager}, apelul \texttt{document.cookie} din Console returneaza \texttt{`session\_id=2322adfdee9be8ef6f73ea2e0e6a2353'} -- proba directa ca lipseste flag-ul \textbf{HttpOnly}, deci un payload XSS reflectat ar putea exfiltra cookie-ul si l-ar putea folosi 30 de zile (max\_age implicit)}{poc-45}

\captura{sessionCookie.png}{Sesiune ``zombie'' dupa logout: dupa ce utilizatorul apasa \emph{Logout} in browser, replay-ul cu \texttt{curl -b "session\_id=\textit{capturat}" /dashboard} returneaza \textbf{HTTP 200} cu pagina logata ca \texttt{manager@deskly.local} -- rândul din tabela \texttt{sessions} nu este sters server-side, deci cookie-ul ramane valid; combinata cu lipsa HttpOnly, vulnerabilitatea permite persistenta atacatorului 30 de zile}{poc-45b}

\subsection{4.6 -- Reset token forjat}
Tokenul fiind \texttt{base64url(user\_id:timestamp)}, atacatorul il
poate genera offline:
\begin{lstlisting}[style=sh]
$ python3 -c "import base64; print(base64.urlsafe_b64encode(b'1:0').decode().rstrip('='))"
MTow
$ # browser: GET /reset/MTow ; submit parola noua "pwn3d-by-attacker"
$ # browser: POST /login cu manager@deskly.local + noua parola -> 302 /dashboard
\end{lstlisting}
Atacatorul preia complet contul managerului fara sa fi avut vreodata
acces la inboxul lui.

\captura{tokenForjat.png}{Pasul 1 al atacului: atacatorul, fara sa fie logat, navigheaza la \texttt{http://127.0.0.1:5000/reset/MTow} (token-ul \texttt{MTow} = \texttt{base64url(b'1:0')}, generat offline cu o comanda \texttt{python3 -c "..."}); aplicatia accepta tokenul si afiseaza formularul \emph{Reset password}, permitand setarea unei noi parole pentru \texttt{user\_id=1}}{poc-46a}

\captura{managerDupatoken.png}{Pasul 2 -- account takeover finalizat: dupa schimbarea parolei via tokenul forjat, atacatorul se logheaza cu \texttt{manager@deskly.local} si noua parola si ajunge pe \texttt{/dashboard} cu rolul \texttt{MANAGER}, fara sa fi avut vreodata acces la inboxul victimei}{poc-46b}

\captura{tokendecodat.png}{Confirmarea structurii tokenurilor reale emise de v1: o comanda Python \texttt{base64.urlsafe\_b64decode} aplicata pe un token capturat de pe \texttt{/forgot} returneaza \texttt{b'5:1777107379'} -- adica \texttt{user\_id:timestamp\_unix}; ambele campuri sunt usor de ghicit (id-uri mici secventiale, timestamp-ul curent)}{poc-46c}

\subsection{4.7 -- SQL injection pe \texttt{/ticket/search}}
Loginat ca \texttt{analyst}, urmatoarele payload-uri demonstreaza
exploatabilitatea:
\begin{lstlisting}[style=sh]
GET /ticket/search?q=
GET /ticket/search?q=
\end{lstlisting}
Prima injectie (\texttt{' OR 1=1--}) inchide ghilimeaua de la
\texttt{LIKE '\%...\%'} si transforma clauza \texttt{WHERE} intr-o
tautologie, iar \texttt{--} comenteaza restul query-ului. Rezultatul:
analist-ul vede \emph{toate} ticketele, inclusiv ticketul confidential
\texttt{\#1 -- Internal report} al managerului. A doua injectie
(UNION-based) expune hash-urile MD5 ale tuturor utilizatorilor in
coloana \texttt{description} a tabelului afisat -- combinabil cu
atacul 4.2 pentru full account takeover.

\captura{ORSQL.png}{SQL injection \texttt{' OR 1=1--}: utilizatorul \texttt{analyst@deskly.local} obtine ambele tickete (inclusiv \texttt{Internal report} cu \texttt{owner\_id=1}, al managerului) -- bara de URL contine payload-ul URL-encoded}{poc-47a}

\captura{sqlUnion.png}{UNION-based SQLi pe \texttt{/ticket/search}: payload-ul \texttt{' UNION SELECT id, email, password\_hash, role, created\_at, locked FROM users--} concateneaza rezultatele cautarii cu intregul tabel \texttt{users}; in coloana \texttt{Description} apar hash-urile MD5 ale tuturor conturilor (\texttt{manager@deskly.local}, \texttt{analyst@deskly.local}, plus eventuale conturi de test) -- combinabil cu atacul 4.2 pentru full account takeover}{poc-47b}

\captura{sqlError.png}{Information disclosure prin mesaje de eroare: un payload UNION trunchiat (\texttt{' UNION SELECT id, email, password\_hash, role, created\_at, locked FROM users-}, fara comentariul final) provoaca un raspuns \emph{``SQL error: near \"-\": syntax error''} afisat direct in pagina; mesajul confirma atacatorului ca engine-ul este SQLite si ca injectia este exploatabila, accelerand tunarea payload-ului}{poc-47c}

\subsection{4.8 -- IDOR pe \texttt{/ticket/<id>}}
Loginat ca \texttt{analyst} (\texttt{user.id=2}), accesarea
\texttt{/ticket/1} (ticketul managerului) returneaza HTTP 200 cu
continutul confidential:
\begin{lstlisting}[style=sh]
GET /ticket/1
=> HTTP/1.1 200 OK
=> <h1>Ticket #1: Internal report</h1>
=> <p>A leaked admin password was found.</p>
\end{lstlisting}

\captura{ticket2.png}{Baseline -- listarea \texttt{/tickets} ca \texttt{analyst@deskly.local}: aplicatia afiseaza un singur rand (\texttt{Phishing email}, \texttt{owner\_id=2}), respectand izolarea legitima a datelor pe owner; lipsa verificarii \texttt{owner\_id} apare doar pe \texttt{/ticket/<id>}}{poc-48a}

\captura{idor.png}{IDOR pur: acelasi \texttt{analyst@deskly.local} (badge sus-dreapta) navigheaza direct la \texttt{/ticket/1} si primeste \texttt{Ticket \#1 -- Internal report}, severitate \texttt{HIGH}, descriere \emph{``A leaked admin password was found.''}; \texttt{owner\_id=1} (manager) este afisat in pagina, iar HTTP 200 in Network tab confirma absenta oricarui control de acces}{poc-48b}

\subsection{4.9 -- Lipsa audit logging}
Dupa ce am executat toate cele 8 atacuri de mai sus, tabela
\texttt{audit\_logs} este goala:
\begin{lstlisting}[style=sh]
$ sqlite3 deskly_vuln.db "SELECT COUNT(*) FROM audit_logs;"
0
\end{lstlisting}

\captura{auditzero.png}{Dovada lipsei audit-ului in v1: \texttt{SELECT COUNT(*) FROM audit\_logs} returneaza \texttt{0} chiar si dupa cele 8 atacuri precedente, iar \texttt{.schema audit\_logs} confirma ca tabela exista in schema seed dar nu este folosita de \texttt{vulnerable/app.py}}{poc-49}

\newpage
\section{Analiza impactului}\label{sec:impact}

Impactul fiecarei vulnerabilitati este evaluat pe triada CIA
(\textit{Confidentiality}, \textit{Integrity}, \textit{Availability})
si exprimat in scenarii concrete de business.

\begin{longtable}{p{0.04\textwidth} p{0.22\textwidth} p{0.10\textwidth} p{0.50\textwidth}}
\toprule
\textbf{ID} & \textbf{Vuln.} & \textbf{CIA} & \textbf{Scenariu de business} \\
\midrule
4.1 & Password policy & C & Parolele triviale permit credential stuffing pe scara larga; conturile reale ale angajatilor sunt compromise prin liste publice (ex. RockYou). \\
4.2 & Stocare MD5 & C & Un dump al DB (backup, exfiltrare via SQLi) expune toate parolele; intrucat utilizatorii reutilizeaza parole, conturile lor sunt compromise si pe alte servicii. \\
4.3 & Brute-force & C, I & Atacatorul preia conturi prin dictionary attack online; nu exista lockout, deci atacul este invizibil pentru SOC. \\
4.4 & User enumeration & C & Atacatorul construieste o lista de email-uri valide, pe care le foloseste tintit la phishing si spear-phishing. \\
4.5 & Sesiune nesigura & C, I & XSS reflectat poate exfiltra cookie-ul (lipsa HttpOnly); cookie-ul capturat este valid 30 de zile si nu este invalidat la logout. \\
4.6 & Reset token & C, I, A & Account takeover total fara interactiune cu victima; atacatorul poate bloca accesul utilizatorului prin schimbarea parolei. \\
4.7 & SQL injection & C, I, A & Citire tabele arbitrare (\texttt{users}, \texttt{audit\_logs}); pe SQLite se poate face si \texttt{ATTACH DATABASE} (write); pierderi de date masive. \\
4.8 & IDOR & C, I & Acces cross-tenant la informatii confidentiale (in scenariul de business: rapoarte de audit intern, tichete cu credentiale). \\
4.9 & Lipsa audit & I, A & Incidentele nu pot fi reconstituite forensic; non-repudierea este imposibila; reglementarile (GDPR Art. 30, ISO 27001 A.12.4) nu sunt respectate. \\
\bottomrule
\caption{Analiza impactului per vulnerabilitate}
\label{tab:impact}
\end{longtable}

\subsection{Scenariu agregat de exploatare}
Cele 9 vulnerabilitati pot fi inlantuite intr-un \textit{kill chain}
realist:
\begin{enumerate}
  \item \textbf{4.4 + 4.3}: atacatorul enumera email-urile valide si
        lanseaza un dictionary attack online; un cont cu parola slaba
        cade in cateva minute.
  \item \textbf{4.7}: cu un cont de \texttt{ANALYST}, atacatorul exfiltreaza
        intreaga tabela \texttt{users} prin SQL injection.
  \item \textbf{4.2}: hash-urile MD5 sunt sparte offline aproape instant
        pentru orice parola din top 10.000.
  \item \textbf{4.6} (alternativ, daca o parola nu cade): atacatorul
        forjeaza tokenul de reset al contului \texttt{MANAGER} si il
        preia direct.
  \item \textbf{4.8}: cu un cont de manager, atacatorul iese complet
        din silosul de vizibilitate si citeste toate ticketele.
  \item \textbf{4.5}: pentru persistenta, atacatorul exfiltreaza
        cookie-ul printr-un XSS si il pastreaza valid 30 de zile.
  \item \textbf{4.9}: pentru ca nicio actiune nu este auditata,
        compromiterea poate trece nedetectata pana la urmatorul audit
        extern.
\end{enumerate}
Concluzie: scorul agregat al riscului este \textbf{Critic}; v1 nu poate
fi expusa in mediu de productie.

\newpage
\section{Implementarea fix-urilor}\label{sec:fix}

Aceasta sectiune prezinta, pentru fiecare vulnerabilitate, codul
modificat in v2 si justificarea controlului ales. Numerele de linie
fac referire la fisierul \texttt{fixed/app.py}.

\subsection{4.1 -- Politica de parole}
Validarea se face prin expresie regulata, aplicata atat la \texttt{Register}
cat si la \texttt{Reset}.

\begin{lstlisting}[style=py, caption={Validare parola in v2}]
PASSWORD_POLICY_RE = re.compile(
    r"^(?=.*[A-Z])(?=.*\d)(?=.*[^A-Za-z0-9]).{8,}$"
)
def validate_password(pw: str) -> str | None:
    if not PASSWORD_POLICY_RE.match(pw):
        return ("Password must be at least 8 chars and include "
                "uppercase, digit and symbol.")
    return None
\end{lstlisting}

\subsection{4.2 -- Hash bcrypt cost 12}
\begin{lstlisting}[style=py, caption={Hash si verificare bcrypt}]
import bcrypt
def hash_password(pw: str) -> str:
    return bcrypt.hashpw(pw.encode(), bcrypt.gensalt(rounds=12)).decode()
def verify_password(pw: str, h: str) -> bool:
    try:
        return bcrypt.checkpw(pw.encode(), h.encode())
    except ValueError:
        return False
\end{lstlisting}
\textbf{Justificare.} Cost factor 12 ofera $\sim 250$~ms per verificare
pe hardware-ul VM, ceea ce face brute-force-ul offline impracticabil.
Salt-ul este generat per parola, deci doua conturi cu aceeasi parola au
hash-uri diferite.

\subsection{4.3 -- Rate limiting si lockout}
Tabela \texttt{users} primeste doua coloane suplimentare:
\texttt{failed\_attempts} si \texttt{lock\_until}. Dupa 5 esecuri
consecutive, contul este blocat 15 minute.

\begin{lstlisting}[style=py, caption={Lockout per cont in v2}]
MAX_FAILED = 5
LOCK_MINUTES = 15

def login_check(user, password) -> tuple[bool, str | None]:
    if user["lock_until"] and user["lock_until"] > _sql_ts():
        return False, "Account locked. Try again later."
    if not verify_password(password, user["password_hash"]):
        n = user["failed_attempts"] + 1
        if n >= MAX_FAILED:
            db.execute("UPDATE users SET failed_attempts=0, "
                       "lock_until=? WHERE id=?",
                       (_sql_ts(minutes=LOCK_MINUTES), user["id"]))
        else:
            db.execute("UPDATE users SET failed_attempts=? WHERE id=?",
                       (n, user["id"]))
        db.commit()
        return False, "Invalid credentials"
    db.execute("UPDATE users SET failed_attempts=0 WHERE id=?",
               (user["id"],))
    return True, None
\end{lstlisting}

\subsection{4.4 -- Mesaj generic si timp uniform}
Pentru a egaliza timpul de raspuns intre ramurile ``user gasit'' si
``user lipsa'', in cazul lipsei se calculeaza un bcrypt dummy.

\begin{lstlisting}[style=py, caption={Mesaj generic + dummy bcrypt}]
DUMMY_HASH = bcrypt.hashpw(b"dummy", bcrypt.gensalt(rounds=12)).decode()
GENERIC = "Invalid credentials"

def login_route():
    user = db.execute("SELECT * FROM users WHERE email=?", (email,)).fetchone()
    if user is None:
        bcrypt.checkpw(b"dummy", DUMMY_HASH.encode())   # timing
        audit("LOGIN_FAIL", resource="auth", resource_id=email)
        return render_template("login.html", error=GENERIC)
    ok, _ = login_check(user, password)
    if not ok:
        audit("LOGIN_FAIL", user_id=user["id"], resource="auth")
        return render_template("login.html", error=GENERIC)
    ...
\end{lstlisting}

\subsection{4.5 -- Cookie hardening si rotatia sesiunii}
\begin{lstlisting}[style=py, caption={Cookie hardening + revoke server-side}]
resp.set_cookie(
    SESSION_COOKIE, token,
    httponly=True, samesite="Lax", secure=False,  # secure=True in HTTPS
    max_age=3600,                                  # 1h, nu 30 zile
)
# La logout:
db.execute("UPDATE sessions SET revoked=1 WHERE id=?", (token,))
db.commit()
\end{lstlisting}
Rotatia sesiunii: la fiecare login se genereaza un \textbf{nou}
\texttt{session\_id}, iar valoarea anterioara (daca exista) este
marcata \texttt{revoked=1}. Astfel, fixarea sesiunii este imposibila.

\subsection{4.6 -- Token aleator, hash-uit, one-time, cu expirare}
\begin{lstlisting}[style=py, caption={Reset token securizat}]
import secrets, hashlib

def make_reset_token(user_id: int) -> str:
    raw = secrets.token_urlsafe(32)
    h = hashlib.sha256(raw.encode()).hexdigest()
    db.execute(
        "INSERT INTO password_resets (user_id, token_hash, "
        "expires_at, used) VALUES (?, ?, ?, 0)",
        (user_id, h, _sql_ts(hours=1)),
    )
    db.commit()
    return raw     # se trimite prin email; in laborator se afiseaza in pagina

def consume_reset_token(raw: str) -> int | None:
    h = hashlib.sha256(raw.encode()).hexdigest()
    row = db.execute(
        "SELECT * FROM password_resets WHERE token_hash=? AND used=0 "
        "AND expires_at > ?", (h, _sql_ts())).fetchone()
    if row is None:
        return None
    db.execute("UPDATE password_resets SET used=1 WHERE id=?", (row["id"],))
    db.commit()
    return row["user_id"]
\end{lstlisting}

\subsection{4.7 -- Query parametrizate}
\begin{lstlisting}[style=py, caption={Search securizat}]
@app.route("/ticket/search")
@login_required
def ticket_search():
    q = request.args.get("q", "").strip()[:100]   # input validation
    like = f"
    rows = get_db().execute(
        "SELECT id, title, description, severity, status, owner_id "
        "FROM tickets WHERE (title LIKE ? OR description LIKE ?) "
        "  AND (owner_id = ? OR ? = 'MANAGER') "
        "ORDER BY id DESC LIMIT 50",
        (like, like, current_user()["id"], current_user()["role"]),
    ).fetchall()
    return render_template("ticket_search.html", q=q, tickets=rows)
\end{lstlisting}

\subsection{4.8 -- Verificare ownership}
\begin{lstlisting}[style=py, caption={IDOR prevention pe \texttt{/ticket/<id>}}]
@app.route("/ticket/<int:ticket_id>")
@login_required
def ticket_detail(ticket_id):
    user = current_user()
    row = get_db().execute(
        "SELECT * FROM tickets WHERE id=?", (ticket_id,)).fetchone()
    if row is None:
        abort(404)
    if row["owner_id"] != user["id"] and user["role"] != "MANAGER":
        audit("IDOR_BLOCKED", user_id=user["id"],
              resource="ticket", resource_id=str(ticket_id))
        abort(403)
    return render_template("ticket_detail.html", ticket=row)
\end{lstlisting}

\subsection{4.9 -- Audit logging}
Helper-ul \texttt{audit()} este apelat in toate ramurile sensibile:
login (succes/esec), reset (cerut/folosit), IDOR blocat, CSRF blocat,
schimbare ticket etc. Tabela \texttt{audit\_logs} este consultabila
prin endpoint-ul \texttt{/audit} (doar pentru rolul \texttt{MANAGER}).

\begin{lstlisting}[style=py, caption={Helper de audit}]
def audit(action: str, *, user_id: int | None = None,
          resource: str = "", resource_id: str = ""):
    get_db().execute(
        "INSERT INTO audit_logs (user_id, action, resource, "
        "resource_id, ip_address) VALUES (?, ?, ?, ?, ?)",
        (user_id, action, resource, resource_id, request.remote_addr),
    )
    get_db().commit()
\end{lstlisting}

\subsection{Controale orizontale}
Pe langa fix-urile dedicate fiecarei vulnerabilitati, in v2 am adaugat
controale orizontale care imbunatatesc postura globala de securitate:

\begin{itemize}
  \item \textbf{CSRF}. Token aleator per sesiune, validat in
        \texttt{before\_request} pentru toate metodele non-safe (POST,
        PUT, DELETE).
  \item \textbf{Security headers}. \texttt{after\_request} adauga
        \texttt{X-Content-Type-Options: nosniff},
        \texttt{X-Frame-Options: DENY},
        \texttt{Referrer-Policy: same-origin}.
  \item \textbf{Content Security Policy}. Directiva
        \texttt{default-src 'self'; script-src 'self'} -- previne
        injectia inline de scripturi (bonus, sectiunea \ref{sec:bonus}).
  \item \textbf{Error handling generic}. Sablonul \texttt{error.html}
        afiseaza un mesaj scurt; stack-trace-ul nu mai ajunge la client.
\end{itemize}

\newpage
\section{Re-test (validarea fix-urilor)}\label{sec:retest}

Re-testul foloseste exact aceleasi PoC-uri ca in capitolul \ref{sec:poc},
dar lansate impotriva versiunii \texttt{fixed}. Rezultatul fiecarui PoC
este consemnat in tabelul \ref{tab:retest} si insotit de captura
corespunzatoare.

\begin{longtable}{p{0.04\textwidth} p{0.32\textwidth} p{0.30\textwidth} p{0.22\textwidth}}
\toprule
\textbf{ID} & \textbf{PoC reluat} & \textbf{Comportament v2} & \textbf{Captura} \\
\midrule
\endhead
4.1 & Register cu \texttt{password=a} & Eroare: ``Password must be at least 8 chars, uppercase, digit, symbol.'' & \ref{fig:retest-41} \\
4.2 & Citire \texttt{password\_hash} din DB & Hash-uri \texttt{\$2b\$12\$...} (bcrypt) & \ref{fig:retest-42} \\
4.3 & 6 incercari gresite & La a 5-a: ``Account locked. Try again later.'' & \ref{fig:retest-43} \\
4.4 & User valid vs.\ inexistent & Acelasi mesaj ``Invalid credentials''; timp uniform & \ref{fig:retest-44} \\
4.5 & Inspectie cookie + replay post-logout & HttpOnly + SameSite + rotatie; replay returneaza 302 \texttt{/login} & \ref{fig:retest-45} \\
4.6 & Token forjat \texttt{MTow} & ``Invalid or expired token'' & \ref{fig:retest-46} \\
4.7 & \texttt{?q=' OR 1=1--} & 0 rezultate (parametru tratat literal) & \ref{fig:retest-47} \\
4.8 & \texttt{/ticket/1} ca analyst & HTTP 403; \texttt{IDOR\_BLOCKED} in audit & \ref{fig:retest-48} \\
4.9 & \texttt{COUNT(*) FROM audit\_logs} & $> 0$, multiple actiuni distincte & \ref{fig:retest-49} \\
\bottomrule
\caption{Sumarizarea re-testului pe v2}
\label{tab:retest}
\end{longtable}

\captura{registerFix.png}{Re-test 4.1: parola \texttt{a} este respinsa la \texttt{/register} cu mesajul \emph{``Password must be at least 8 characters and contain an uppercase letter, a digit, and a special character''}; \texttt{PASSWORD\_POLICY\_RE} se aplica si la \texttt{/reset}}{retest-41}

\captura{parolaBcrypt.png}{Re-test 4.2: dupa migrare, fiecare \texttt{password\_hash} din \texttt{deskly\_fixed.db} incepe cu prefixul \texttt{\$2b\$12\$} (bcrypt v2b, cost 12, lungime 60 caractere) -- nu mai poate fi cracat offline cu rainbow tables}{retest-42}

\captura{loginlock.png}{Re-test 4.3: dupa 5 esecuri consecutive contul \texttt{locktest\_*} este blocat 15 minute; in DB \texttt{failed\_attempts=5} si \texttt{lock\_until} este timestamp in viitor; chiar si parola corecta nu mai trece pana la expirare}{retest-43}

\captura{invalidCredentials.png}{Re-test 4.4 (UI): pentru orice combinatie email + parola gresita aplicatia afiseaza acelasi mesaj generic \emph{``Invalid credentials''}, fara sa deconspire daca emailul exista in sistem}{retest-44a}

\captura{invalidCredCurl.png}{Re-test 4.4 (raspuns HTTP): \texttt{curl} pe \texttt{/login} returneaza acelasi corp HTML pentru user inexistent si pentru parola gresita; combinat cu hash-ul dummy bcrypt aplicat pe ramura ``user not found'', timpul de raspuns ramane uniform}{retest-44b}

\captura{noIdHttpOnly.png}{Re-test 4.5 (HttpOnly): in browser-ul logat ca manager, \texttt{document.cookie} returneaza string vid -- cookie-ul \texttt{session\_id} nu mai poate fi citit din JavaScript, deci un payload XSS nu il mai poate exfiltra}{retest-45a}

\captura{cookie1.png}{Re-test 4.5 (DevTools $\rightarrow$ Cookies): cookie-ul \texttt{session\_id} are flag-urile \textbf{HttpOnly} si \textbf{SameSite=Lax} active si o expirare scurta (\texttt{max\_age=3600}), spre deosebire de v1 unde ramanea valid 30 de zile}{retest-45b}

\captura{cookie2.png}{Re-test 4.5 (rotatie): dupa logout + re-login, valoarea cookie-ului se schimba complet -- noul \texttt{session\_id} difera de cel anterior, blocand atacurile de tip session fixation}{retest-45c}

\captura{curlCookie.png}{Re-test 4.5 (revoke server-side): replay cu \texttt{curl -b "session\_id=\textit{vechi}"} pe \texttt{/dashboard} dupa logout primeste status \textbf{302} catre \texttt{/login} -- sesiunea a fost invalidata in tabela \texttt{sessions} (coloana \texttt{revoked=1}), nu mai e suficient sa stergi cookie-ul client-side}{retest-45d}

\captura{invalidExpiredTokenMtow.png}{Re-test 4.6: acelasi payload \texttt{/reset/MTow} care a permis takeover-ul in v1 este respins de v2 cu mesajul \emph{``Invalid or expired token''}; tokenul real este acum 32 bytes random, stocat hash-uit (SHA-256), expira in 1h si este single-use}{retest-46}

\captura{resetpassword.png}{Re-test 4.6 (continuare): chiar si pe un token valid, formularul \texttt{/reset} aplica aceeasi politica de parole ca \texttt{/register} -- o parola noua \texttt{`pwn3d-by-attacker'} ar fi respinsa de \texttt{PASSWORD\_POLICY\_RE}}{retest-46b}

\captura{tokenuriGenerate.png}{Re-test 4.6 (DB schema noua): tabela \texttt{password\_resets} din \texttt{deskly\_fixed.db} stocheaza in coloana \texttt{token\_hash} doar SHA-256 al tokenului real (siruri de 64 caractere hex), nu valoarea brută; coloanele \texttt{expires\_at} (1 ora dupa emitere) si \texttt{used} (0/1) implementeaza expirare si single-use; chiar si un atacator cu acces la DB nu poate folosi tokenurile inca emise}{retest-46c}

\captura{tokenCantBeUsedTwice.png}{Re-test 4.6 (single-use): dupa folosirea unui token valid pentru schimbarea parolei, o a doua incercare cu acelasi token este respinsa cu acelasi mesaj generic \emph{``Invalid or expired token''}; in DB, randul corespunzator are \texttt{used=1} (vizibil in panoul stang), iar mesajul nu deconspira daca tokenul a fost vreodata valid}{retest-46d}

\captura{sqlInjection.png}{Re-test 4.7 (\texttt{OR 1=1}): acelasi payload \texttt{' OR 1=1--} care a expus toate ticketele in v1 produce acum \emph{``No results for ' OR 1=1--''} -- aplicatia afiseaza explicit banner-ul verde \emph{``Queries are bound as parameters -- SQL injection payloads are treated as literal text''}, confirmand ca input-ul este pasat ca parametru \texttt{?} catre SQLite}{retest-47a}

\captura{sql2Union.png}{Re-test 4.7 (UNION): si payload-ul mai agresiv \texttt{' UNION SELECT id, email, password\_hash, role, created\_at, locked FROM users--} este neutralizat -- 0 rezultate, niciun \texttt{SQL error: ...} expus, niciun hash leak-uit; coloana \texttt{Description} ramane tratata ca text}{retest-47b}

\captura{audit.png}{Re-test 4.8: navigarea la o resursa pentru care utilizatorul nu are drepturi (in cazul de fata \texttt{/audit} ca \texttt{analyst}) este blocata cu \textbf{HTTP 403 Forbidden}; acelasi mecanism (\texttt{owner\_id} sau rol \texttt{MANAGER}) se aplica la \texttt{/ticket/<id>}, iar tentativa este logata ca \texttt{IDOR\_BLOCKED} / \texttt{AUTHZ\_DENY} in audit (vizibil in panoul stang)}{retest-48}

\captura{auditSqlAggregation.png}{Re-test 4.9: \texttt{SELECT action, COUNT(*) FROM audit\_logs GROUP BY action} pe \texttt{deskly\_fixed.db} returneaza zeci de inregistrari distribuite pe \texttt{LOGIN\_SUCCESS}, \texttt{LOGIN\_FAIL}, \texttt{LOGOUT}, \texttt{REGISTER}, \texttt{RESET\_REQUESTED}, \texttt{RESET\_USED}, \texttt{IDOR\_BLOCKED}, \texttt{ACCOUNT\_LOCKED} -- vs.\ 0 in v1}{retest-49}

\subsection{Test de regresie automatizat}
Pe langa re-testul manual, am inclus doua scripturi de smoke:
\texttt{smoke\_vuln.py} ruleaza toate cele 9 PoC-uri pe v1 si confirma
ca atacurile reusesc; \texttt{smoke\_fixed.py} ruleaza aceleasi PoC-uri
pe v2 si confirma ca fiecare este blocat. Ambele scripturi se incheie
cu \texttt{All vulnerable PoCs reproduced} respectiv \texttt{All fixed
PoCs hold up under attack}, fara nicio asertie esuata. Acest mecanism
acopera bonusul ``Teste automate de securitate'' (sectiunea \ref{sec:bonus}).

\begin{lstlisting}[style=sh, caption={Rularea scripturilor de smoke}]
$ python smoke_vuln.py     # cu v1 pornit
[4.1] weak password accepted on register
  PASS -- register with single-char password reached the server
  ... (toate cele 9)
All vulnerable PoCs reproduced.
$ python smoke_fixed.py    # cu v2 pornit
[4.1] register rejects weak passwords
  PASS -- single-char password rejected
  ... (toate cele 9)
All fixed PoCs hold up under attack.
\end{lstlisting}

\newpage
\section{Audit \& logging}\label{sec:audit}

\subsection{Modelul de audit}
In v2, fiecare actiune sensibila inregistreaza un rand in tabela
\texttt{audit\_logs}, cu campurile prezentate in tabelul \ref{tab:audit}.

\begin{table}[H]
\centering
\begin{tabular}{ll}
\toprule
\textbf{Camp} & \textbf{Continut} \\
\midrule
\texttt{user\_id}    & ID-ul utilizatorului care a executat actiunea (sau NULL pentru actiuni anonime) \\
\texttt{action}      & Tipul actiunii (\texttt{LOGIN\_SUCCESS}, \texttt{LOGIN\_FAIL}, \texttt{IDOR\_BLOCKED}, ...) \\
\texttt{resource}    & Tipul resursei (\texttt{auth}, \texttt{ticket}, \texttt{audit}) \\
\texttt{resource\_id}& Identificatorul resursei (de ex.\ \texttt{ticket.id}) \\
\texttt{timestamp}   & Momentul actiunii \\
\texttt{ip\_address} & IP-ul clientului \\
\bottomrule
\end{tabular}
\caption{Schema tabelei \texttt{audit\_logs}}
\label{tab:audit}
\end{table}

\subsection{Actiuni inregistrate}
Helper-ul \texttt{audit()} este apelat in:
\begin{itemize}
  \item \textbf{Auth}: \texttt{LOGIN\_SUCCESS}, \texttt{LOGIN\_FAIL},
        \texttt{LOGOUT}, \texttt{REGISTER}, \texttt{ACCOUNT\_LOCKED}.
  \item \textbf{Reset}: \texttt{RESET\_REQUESTED}, \texttt{RESET\_USED},
        \texttt{RESET\_INVALID}.
  \item \textbf{Tickets}: \texttt{TICKET\_CREATED}, \texttt{TICKET\_VIEWED},
        \texttt{IDOR\_BLOCKED}.
  \item \textbf{Securitate}: \texttt{CSRF\_FAIL}.
\end{itemize}

\subsection{Vizualizarea audit log-urilor}
Endpoint-ul \texttt{/audit} este accesibil doar pentru rolul
\texttt{MANAGER} si afiseaza ultimele 200 inregistrari, ordonate
descrescator dupa \texttt{timestamp}.

\captura{auditscreen.png}{Pagina \texttt{/audit} (accesibila doar rolului \texttt{MANAGER}): tabel cronologic cu evenimente \texttt{LOGIN\_SUCCESS}, \texttt{LOGOUT}, \texttt{AUTHZ\_DENY}, \texttt{VIEW\_TICKET}, \texttt{SEARCH}, \texttt{PASSWORD\_RESET}, fiecare cu \texttt{user\_id}, \texttt{resource}, \texttt{ip\_address} si timestamp}{audit-view}

\subsection{Trasabilitate end-to-end}
Demonstratia de trasabilitate consta in:
(i) lansarea unui atac (de ex. brute-force pe contul managerului),
(ii) deschiderea \texttt{/audit} si identificarea celor 5 inregistrari
\texttt{LOGIN\_FAIL} cu acelasi \texttt{ip\_address}, urmate de un
\texttt{ACCOUNT\_LOCKED}. Inregistrarile pot fi exportate (CSV) ca
proba pentru un raspuns la incident.

\begin{lstlisting}[style=sh, caption={Audit dupa o tentativa de brute-force}]
$ sqlite3 deskly_fixed.db \
  "SELECT timestamp, action, resource_id FROM audit_logs ORDER BY id DESC LIMIT 6;"
2026-04-29 11:55:18 | ACCOUNT_LOCKED | manager@deskly.local
2026-04-29 11:55:17 | LOGIN_FAIL     | manager@deskly.local
2026-04-29 11:55:17 | LOGIN_FAIL     | manager@deskly.local
2026-04-29 11:55:16 | LOGIN_FAIL     | manager@deskly.local
2026-04-29 11:55:16 | LOGIN_FAIL     | manager@deskly.local
2026-04-29 11:55:15 | LOGIN_FAIL     | manager@deskly.local
\end{lstlisting}

\newpage
\section{Concluzii}\label{sec:concluzii}

\subsection{Lectii invatate}
\begin{itemize}
  \item \textbf{Securitatea nu este un strat optional.} Fiecare ramura
        de cod care atinge un secret (parola, token, sesiune) trebuie
        gandita din perspectiva atacatorului inca de la prima
        implementare. Adaugarea ulterioara este, intotdeauna, mai
        scumpa si mai fragila.
  \item \textbf{Defense in depth.} Politica de parole (4.1) ar fi
        insuficienta fara hash bcrypt (4.2) si fara rate limiting (4.3).
        Trei controale complementare reduc impactul reciproc daca unul
        cedeaza.
  \item \textbf{Mesajele de eroare sunt suprafata de atac.} Diferenta
        intre ``user does not exist'' si ``incorrect password'' (4.4)
        si scurgerea \texttt{SQL error: ...} (4.7) demonstreaza ca un
        text aparent inofensiv poate accelera dramatic exploatarea.
  \item \textbf{Audit-ul este ne-negociabil.} Fara
        \texttt{audit\_logs} (4.9), atacurile demonstrate aici trec
        nedetectate; in productie, lipsa trasabilitatii incalca
        cerintele de compliance (GDPR Art. 30, ISO 27001).
  \item \textbf{Cost factor bcrypt.} Alegerea cost factor-ului 12
        este un compromis intre rezistenta si latenta. Pe acelasi VM,
        cost 14 a urcat latenta peste 1~s la login -- inutilizabil. Pe
        un server cu CPU mai bun se poate incerca 13 sau 14.
  \item \textbf{Smoke tests automate.} Cele doua scripturi
        (\texttt{smoke\_vuln.py}, \texttt{smoke\_fixed.py}) au prins
        regresii in timpul iteratiei (de ex. o problema de comparare
        lexicografica a timestamp-urilor in tabela \texttt{sessions}).
\end{itemize}

\newpage
\section{Criterii de excelenta (bonus)}\label{sec:bonus}

\subsection{Implementare CSP (+5p)}
Versiunea securizata adauga in fiecare raspuns headerul
\texttt{Content-Security-Policy}, configurat sa permita doar resurse
proprii (\texttt{'self'}). Politica blocheaza inline-script si
inline-style, prevenind clase intregi de XSS reflectat.

\begin{lstlisting}[style=py, caption={CSP in \texttt{after\_request}}]
@app.after_request
def security_headers(resp):
    resp.headers["Content-Security-Policy"] = (
        "default-src 'self'; "
        "script-src 'self'; "
        "style-src 'self'; "
        "img-src 'self' data:; "
        "object-src 'none'; "
        "base-uri 'self'; "
        "frame-ancestors 'none'"
    )
    resp.headers["X-Content-Type-Options"] = "nosniff"
    resp.headers["X-Frame-Options"] = "DENY"
    resp.headers["Referrer-Policy"] = "same-origin"
    return resp
\end{lstlisting}

\captura{cspheader.png}{Header-ele de securitate emise de v2 la \texttt{GET /login} (\texttt{curl -sI}): \texttt{X-Content-Type-Options: nosniff}, \texttt{X-Frame-Options: DENY}, \texttt{Referrer-Policy: no-referrer} si \texttt{Content-Security-Policy: default-src 'self'; style-src 'self' 'unsafe-inline'; img-src 'self' data:; frame-ancestors 'none'} -- toate absente in v1}{bonus-csp}

\subsubsection*{Protectia CSRF (token per-sesiune)}
Toate metodele non-safe (POST, PUT, DELETE) sunt validate intr-un hook
\texttt{before\_request} care compara campul \texttt{csrf\_token} din
formular cu valoarea stocata in sesiunea server-side. Lipsa sau
nepotrivirea tokenului se finalizeaza cu redirect catre \texttt{/login}
si o intrare \texttt{CSRF\_FAIL} in \texttt{audit\_logs}.

\captura{crsfTokenCreateTicket.png}{Tokenul CSRF este injectat ca \texttt{<input type="hidden" name="csrf\_token">} in fiecare formular; in DevTools $\rightarrow$ Network, body-ul cererii \texttt{POST /ticket/new} contine perechea \texttt{csrf\_token=...} alaturi de campurile utilizatorului}{bonus-csrf-token}

\captura{curlNoCrsf.png}{\texttt{curl -X POST /ticket/new} fara parametrul \texttt{csrf\_token} este blocat de hook-ul \texttt{before\_request} cu un redirect 302; chiar daca atacatorul detine cookie-ul valid de sesiune, nu poate efectua actiuni in numele utilizatorului}{bonus-csrf-curl}

\captura{crsf302.png}{Acelasi POST manual din browser, fara token: pagina e re-incarcata pe \texttt{/login} (status 302), iar Network tab-ul confirma blocarea}{bonus-csrf-302}

\captura{consolaCrsf.png}{Logul Flask: zeci de cereri \texttt{POST /ticket/new HTTP/1.1 302} si \texttt{400} respinse de validatorul CSRF in timpul testului fuzzing -- niciuna nu a produs efecte de scriere in DB}{bonus-csrf-log}

\subsection{Teste automate de securitate (+5p)}
Cele doua scripturi de smoke (\texttt{smoke\_vuln.py},
\texttt{smoke\_fixed.py}) acopera cerinta. Ambele sunt complete, rulate
final si fara asertii esuate. Pot fi integrate trivial in pipeline-ul
CI prin \texttt{python smoke\_fixed.py}.

\end{document}
